\chapter[Band 47: Metrische Räume und Vollständigkeit]{Band 47 auf einen Blick}
\noindent\textbf{Metrische Räume und Vollständigkeit}\par\medskip
\noindent\emph{Lesart.}
Soweit nicht anders angegeben, ist \((X,d)\) ein metrischer Raum,
sind \(A,K,U\subseteq X\), \(x,y,z\in X\), \(a\colon\mathbb N\to X\),
und Radien sowie \(\varepsilon,\delta\) sind positive reelle Zahlen.
Bei Abbildungen \(f\colon X\to Y\) sind \(d_X,d_Y\) die jeweiligen Metriken.
\(\varphi\colon\mathbb N\to\mathbb N\) bezeichnet stets eine strikt
wachsende Teilfolgenabbildung.

\begin{BandUebersicht}
\textbf{Metrik und ihre Axiome}
\(\operatorname{Metrik}_X(d)\) umfasst
\UebFormel{\begin{aligned}
&d\colon X\times X\to\mathbb R,\qquad0\leq d(x,y),\\
&d(x,y)=0\leftrightarrow x=y,\\
&d(x,y)=d(y,x),\\
&d(x,z)\leq d(x,y)+d(y,z).
\end{aligned}}
\UebQuellen{\FormulaRefAuto{MetricSpaceDef}[def];
\FormulaRefAuto{MetricFunctionAxiom}[ax];
\FormulaRefAuto{MetricPositivityAxiom}[ax];
\FormulaRefAuto{MetricDefinitenessAxiom}[ax];
\FormulaRefAuto{MetricSymmetryAxiom}[ax];
\FormulaRefAuto{MetricTriangleAxiom}[ax]}
&
\textit{Umgekehrte Dreiecksungleichung}
\UebFormel{|d(x,z)-d(y,z)|\leq d(x,y).}
Ein Abstand von einem festen Punkt kann sich höchstens um den
Abstand zwischen den beiden bewegten Punkten ändern.
\UebQuellen{\FormulaRefAuto{MetricReverseTriangleInequality}[thm]}
\\ \addlinespace[.8ex]

\textbf{Reelle und diskrete Metrik}
\UebFormel{\begin{aligned}
d_{\mathbb R}(x,y)&\coloneqq|x-y|,\\
\delta_X(x,y)&\coloneqq
\begin{cases}0,&x=y,\\1,&x\neq y.\end{cases}
\end{aligned}}
Die erste Formel gilt für \(x,y\in\mathbb R\), die zweite auf jeder Menge \(X\).
\UebQuellen{\FormulaRefAuto{RealMetricDistanceFunctionDef}[def];
\FormulaRefAuto{DiscreteMetricDef}[def]}
&
\textit{Modelle und diskrete Konvergenz}
\UebFormel{\begin{aligned}
&\operatorname{Metrik}_{\mathbb R}(d_{\mathbb R}),\\
&\operatorname{Metrik}_X(\delta_X),\\
&a_n\MetricTo{\delta_X}x
\leftrightarrow\exists N\;\forall n\geq N\;(a_n=x).
\end{aligned}}
\UebQuellen{\FormulaRefAuto{RealDistanceIsMetric}[thm];
\FormulaRefAuto{DiscreteMetricIsMetric}[thm];
\FormulaRefAuto{DiscreteMetricConvergenceEventuallyConstant}[thm]}
\\ \addlinespace[.8ex]

\textbf{Teilraummetrik}
Für \(Y\subseteq X\):
\UebFormel{d_Y\coloneqq d\restriction_{Y\times Y}.}
\UebQuellen{\FormulaRefAuto{MetricSubspaceRestrictionDef}[def]}
&
\textit{Einschränkung erhält die metrischen Begriffe}
\UebFormel{\operatorname{Metrik}_Y(d_Y).}
Für \(a\colon\mathbb N\to Y\) und \(y\in Y\):
\UebFormel{\begin{aligned}
&a_n\MetricTo{d_Y}y\leftrightarrow a_n\MetricTo d y,\\
&\operatorname{Cauchy}_{d_Y}(a)
\leftrightarrow\operatorname{Cauchy}_d(a).
\end{aligned}}
\UebQuellen{\FormulaRefAuto{MetricSubspaceRestrictionIsMetric}[thm];
\FormulaRefAuto{MetricSubspaceConvergenceAgree}[thm];
\FormulaRefAuto{MetricSubspaceCauchyAgree}[thm]}
\\ \addlinespace[.8ex]

\textbf{Offene Kugeln und offene Mengen}
\UebFormel{\begin{aligned}
&\MetricBall d x r\coloneqq\{y\in X\mid d(y,x)<r\},\\
&\operatorname{Offen}_d(U)\coloneqq\\
&\quad\forall x\in U\;\exists r>0\;
(\MetricBall d x r\subseteq U).
\end{aligned}}
\UebQuellen{\FormulaRefAuto{MetricOpenBallDef}[def];
\FormulaRefAuto{MetricOpenClosedSetDef}[def]}
&
\textit{Kugelverkleinerung und Offenheit}
\UebFormel{\begin{aligned}
&y\in\MetricBall d x r\vdash\exists s>0\\
&\qquad\MetricBall d y s\subseteq\MetricBall d x r,\\
&\operatorname{Offen}_d(\MetricBall d x r).
\end{aligned}}
Für \(d_{\mathbb R}\) sind diese Kugeln die reellen
Epsilon-Umgebungen.
\UebQuellen{\FormulaRefAuto{MetricBallShrinking}[thm];
\FormulaRefAuto{MetricOpenBallsAreOpen}[thm];
\FormulaRefAuto{RealMetricBallIsEpsilonNeighborhood}[thm]}
\\ \addlinespace[.8ex]

\textbf{Metrische Konvergenz}
\UebFormel{\begin{aligned}
a_n\MetricTo d x\coloneqq{}&
\forall\varepsilon>0\;\exists N\in\mathbb N\\
&\forall n\geq N\;d(a_n,x)<\varepsilon.
\end{aligned}}
\UebQuellen{\FormulaRefAuto{MetricSequenceConvergenceDef}[def]}
&
\textit{Eindeutigkeit und Beständigkeit des Grenzwerts}
\UebFormel{\begin{aligned}
&a_n\MetricTo d x,\ a_n\MetricTo d y\vdash x=y,\\
&a_n\MetricTo d x\vdash a_{\varphi(n)}\MetricTo d x.
\end{aligned}}
Endliche Änderungen und Restfolgen erhalten ebenfalls den Grenzwert;
für \(d_{\mathbb R}\) stimmt die Definition mit der reellen Konvergenz überein.
\UebQuellen{\FormulaRefAuto{MetricSequenceLimitUnique}[thm];
\FormulaRefAuto{MetricSubsequencePreservesLimit}[thm];
\FormulaRefAuto{MetricFiniteModificationPreservesLimit}[thm];
\FormulaRefAuto{MetricSequenceTailSameLimit}[thm];
\FormulaRefAuto{RealAndMetricSequenceConvergenceAgree}[thm]}
\\ \addlinespace[.8ex]

\textbf{Beschränkte Mengen und Folgen}
\UebFormel{\begin{aligned}
&\operatorname{Beschr}_d(A)\coloneqq\\
&\quad(A=\varnothing\lor
\exists c\in X\;\exists r>0\\
&\qquad A\subseteq\MetricBall d c r),\\
&\operatorname{Beschr}_d(a)\coloneqq
\operatorname{Beschr}_d(a[\mathbb N]).
\end{aligned}}
\UebQuellen{\FormulaRefAuto{MetricBoundedSetSequenceDef}[def]}
&
\textit{Notwendige, aber keine hinreichende Bedingung}
\UebFormel{a_n\MetricTo d x\vdash\operatorname{Beschr}_d(a).}
Die diskrete Menge \((\mathbb N,\delta_{\mathbb N})\) ist beschränkt,
aber nicht folgenkompakt.
\UebQuellen{\FormulaRefAuto{ConvergentMetricSequenceBounded}[thm];
\FormulaRefAuto{BoundedDiscreteNaturalNumbersNotSequentiallyCompact}[thm]}
\\ \addlinespace[.8ex]

\textbf{Metrische Cauchy-Folge}
\UebFormel{\begin{aligned}
&\operatorname{Cauchy}_d(a)\coloneqq\\
&\quad\forall\varepsilon>0\;\exists N\in\mathbb N\\
&\qquad\forall m,n\geq N\;d(a_m,a_n)<\varepsilon.
\end{aligned}}
\UebQuellen{\FormulaRefAuto{MetricCauchySequenceDef}[def]}
&
\textit{Cauchy-Folgen und konvergente Teilfolgen}
\UebFormel{\begin{aligned}
&a_n\MetricTo d x\vdash\operatorname{Cauchy}_d(a),\\
&\operatorname{Cauchy}_d(a)\vdash\operatorname{Beschr}_d(a),\\
&\operatorname{Cauchy}_d(a),\
a_{\varphi(n)}\MetricTo d x\\
&\qquad\vdash a_n\MetricTo d x.
\end{aligned}}
\UebQuellen{\FormulaRefAuto{MetricConvergentSequenceIsCauchy}[thm];
\FormulaRefAuto{MetricCauchySequenceBounded}[thm];
\FormulaRefAuto{MetricCauchyWithConvergentSubsequence}[thm]}
\\ \addlinespace[.8ex]

\textbf{Vollständigkeit}
\UebFormel{\begin{aligned}
&\operatorname{Vollst}(X,d)\coloneqq\\
&\quad\forall a\colon\mathbb N\to X\\
&\qquad(\operatorname{Cauchy}_d(a)\to\\
&\qquad\exists x\in X\;(a_n\MetricTo d x)).
\end{aligned}}
\UebQuellen{\FormulaRefAuto{CompleteMetricSpaceDef}[def]}
&
\textit{Vollständige und unvollständige Modelle}
\UebFormel{\begin{aligned}
&\operatorname{Vollst}(\mathbb R,d_{\mathbb R}),\\
&\operatorname{Vollst}(X,\delta_X),\\
&\neg\operatorname{Vollst}
(\mathbb R\setminus\{0\},(d_{\mathbb R})_{\mathbb R\setminus\{0\}}).
\end{aligned}}
Die letzte Aussage wird durch eine Cauchy-Folge mit fehlendem
Grenzwert \(0\) im Teilraum belegt.
\UebQuellen{\FormulaRefAuto{RealMetricSpaceComplete}[thm];
\FormulaRefAuto{DiscreteMetricSpaceComplete}[thm];
\FormulaRefAuto{PuncturedRealLineIncomplete}[thm]}
\\ \addlinespace[.8ex]

\textbf{Abgeschlossenheit}
\UebFormel{\begin{aligned}
\operatorname{Abgeschlossen}_d(A)
\coloneqq\operatorname{Offen}_d(X\setminus A).
\end{aligned}}
\UebQuellen{\FormulaRefAuto{MetricOpenClosedSetDef}[def]}
&
\textit{Grenzwerte und vollständige Teilräume}
Ist \(A\) abgeschlossen, so gilt für \(a\colon\mathbb N\to A\):
\UebFormel{a_n\MetricTo d x\vdash x\in A.}
Ist außerdem \(X\) vollständig, dann
\UebFormel{\operatorname{Vollst}(A,d_A).}
\UebQuellen{\FormulaRefAuto{MetricClosedSetContainsSequenceLimits}[thm];
\FormulaRefAuto{ClosedSubspaceOfCompleteMetricSpace}[thm]}
\\ \addlinespace[.8ex]

\textbf{Stetigkeit}
\UebFormel{\begin{aligned}
&\operatorname{Stetig}_{d_X,d_Y}(f,x)\coloneqq\\
&\quad\forall\varepsilon>0\;\exists\delta>0\;\forall u\in X:\\
&\qquad d_X(u,x)<\delta\\
&\qquad\to d_Y(f(u),f(x))<\varepsilon.
\end{aligned}}
Stetigkeit von \(f\) bedeutet Stetigkeit in jedem \(x\in X\).
\UebQuellen{\FormulaRefAuto{MetricContinuityDef}[def]}
&
\textit{Folgenkriterium und Komposition}
\UebFormel{\begin{aligned}
&\operatorname{Stetig}_{d_X,d_Y}(f,x)\\
&\ \leftrightarrow
\forall a\colon\mathbb N\to X:\\
&\qquad(a_n\MetricTo{d_X}x
\to f(a_n)\MetricTo{d_Y}f(x)).
\end{aligned}}
Für stetige \(f\colon X\to Y\) und \(g\colon Y\to Z\)
ist auch \(g\circ f\) stetig.
\UebQuellen{\FormulaRefAuto{MetricSequentialContinuityCriterion}[thm];
\FormulaRefAuto{MetricContinuousComposition}[thm]}
\\ \addlinespace[.8ex]

\textbf{Isometrie}
\UebFormel{\begin{aligned}
&\operatorname{Isometrie}_{d_X,d_Y}(f)\coloneqq\\
&\quad\forall x,u\in X\\
&\qquad d_Y(f(x),f(u))=d_X(x,u).
\end{aligned}}
\UebQuellen{\FormulaRefAuto{MetricIsometryDef}[def]}
&
\textit{Erhalt der metrischen Struktur}
Unter der Isometrieannahme:
\UebFormel{\begin{aligned}
&\operatorname{Stetig}_{d_X,d_Y}(f),\\
&\operatorname{Cauchy}_{d_X}(a)\vdash
\operatorname{Cauchy}_{d_Y}(f\circ a).
\end{aligned}}
\UebQuellen{\FormulaRefAuto{MetricIsometryPreservesStructure}[thm]}
\\ \addlinespace[.8ex]

\textbf{Folgenkompaktheit}
\UebFormel{\begin{aligned}
&\operatorname{Folgenkompakt}_d(K)\coloneqq\\
&\quad\forall a\colon\mathbb N\to K\;\exists x\in K\\
&\quad\exists\varphi\text{ strikt wachsend}:\\
&\qquad a_{\varphi(n)}\MetricTo d x.
\end{aligned}}
Der Grenzwert muss zur betrachteten Teilmenge \(K\) gehören.
\UebQuellen{\FormulaRefAuto{MetricSequentialCompactnessDef}[def]}
&
\textit{Strukturelle Folgerungen}
Jede folgenkompakte Teilmenge ist beschränkt und abgeschlossen;
als metrischer Teilraum ist sie vollständig.
\UebFormel{\operatorname{Folgenkompakt}_d(K)
\vdash\operatorname{Vollst}(K,d_K).}
Endliche Teilmengen sind folgenkompakt, und stetige Bilder
folgenkompakter Mengen sind folgenkompakt.
\UebQuellen{\FormulaRefAuto{MetricSequentiallyCompactSetBounded}[thm];
\FormulaRefAuto{MetricSequentiallyCompactSetClosed}[thm];
\FormulaRefAuto{MetricSequentiallyCompactSetComplete}[thm];
\FormulaRefAuto{FiniteMetricSetSequentiallyCompact}[thm];
\FormulaRefAuto{MetricContinuousImageSequentiallyCompact}[thm]}
\\ \addlinespace[.8ex]

\textbf{Reelle Spezialisierung}
Die metrischen Begriffe werden mit
\UebFormel{X=\mathbb R,\qquad d_{\mathbb R}(x,y)=|x-y|}
verwendet. Eine reelle Folge ist beschränkt, wenn ihre Wertemenge
in einer Kugel liegt.
\UebQuellen{\FormulaRefAuto{RealMetricDistanceFunctionDef}[def];
\FormulaRefAuto{MetricBoundedSetSequenceDef}[def]}
&
\textit{Bolzano--Weierstraß und kompakte Intervalle}
Jede beschränkte reelle Folge besitzt eine konvergente Teilfolge:
\UebFormel{\begin{aligned}
&a\colon\mathbb N\to\mathbb R,\ \operatorname{Beschr}(a)\\
&\qquad\vdash\exists L\in\mathbb R\;\exists\varphi
\text{ strikt wachsend}:\\
&\hspace{7em}a_{\varphi(n)}\to L.
\end{aligned}}
Für reelle \(u\leq v\) ist \([u,v]_{\mathbb R}\) folgenkompakt.
\UebQuellen{\FormulaRefAuto{RealSequenceMonotoneSubsequence}[thm];
\FormulaRefAuto{RealSequenceBolzanoWeierstrass}[thm];
\FormulaRefAuto{RealClosedBoundedIntervalSequentiallyCompact}[thm]}
\\
\end{BandUebersicht}

\noindent\emph{Anschluss.}
Aus einer Norm entsteht der Abstand \(d(x,y)=\lVert x-y\rVert\).
Der abschließende Ausblick ordnet dies in die spätere Theorie der Vektorräume ein.
Die hier entwickelten Begriffe benötigen auf \(X\) selbst keine
Addition oder Skalarmultiplikation.\par
